
% !TeX document-id = {d8b4925c-2057-42a4-b894-2f1a3f1b6345}
%!TeX TXS-program:compile = txs:///xelatex/[--shell-escape]
\documentclass[aspectratio=169, mathserif]{beamer}% TPU recommends 16:9 ratio, 4:3 may require some work with inner theme .sty file

% Style options:
% light --- light theme (default)
% dark --- dark theme
% enlogo --- english TPU logo {default}
% rulogo --- russian TPU logo

\usetheme[light, rulogo]{tpu}% dark theme used as an example of optional argument

\usepackage[russian]{babel}%uncomment this to work in russian
\usepackage[utf8]{inputenc}

\usepackage{fontspec}

\setromanfont{Brygada1918}[
Path=./fonts/BrygadaFontFiles/,
Extension = .ttf,
UprightFont=*-Regular,
BoldFont=*-Bold,
ItalicFont=*-Italic,
BoldItalicFont=*-BoldItalic
]

\setsansfont{ALSSirius}[
Path=./fonts/ALSSiriusFiles/,
Extension = .otf,
UprightFont=*-Regular,
BoldFont=*-Bold,
%ItalicFont=*-Italic,
%BoldItalicFont=*-BoldItalic
]

\setmonofont{Consolas}[
Path=./fonts/ConsolasFontFiles/,
%Scale=0.85,
Extension = .ttf,
UprightFont=*-Regular,
BoldFont=*-Bold,
ItalicFont=*-Italic,
BoldItalicFont=*-BoldItalic
]

\usepackage[cache=false]{minted}
\usepackage{xcolor} % to access the named colour LightGray
\definecolor{LightGray}{gray}{0.9}
\definecolor{onedarkBckGr}{RGB}{40, 44, 52}

\usemintedstyle[python]{default}
\setminted[python]{
fontsize=\scriptsize,
escapeinside=||,
mathescape=true,
numbersep=5pt,
gobble=2,
linenos=false,
frame=single,
framesep=1mm,
python3=true,
bgcolor=backcolour,
}

\usemintedstyle[pycon]{default}
\setminted[pycon]{
	fontsize=\scriptsize,
	escapeinside=||,
	mathescape=true,
	numbersep=5pt,
	gobble=2,
	linenos=false,
	frame=single,
	framesep=1mm,
	python3=true,
	bgcolor=backcolour,
	linenos=true,
}

\newmint{python}{}

\usepackage{booktabs}% good looking tables
\usepackage{multicol}% text in multiple columns, useful for side-by-side text and pictures
\usepackage{hyperref}
\definecolor{maroon}{cmyk}{0, 0.87, 0.68, 0.32}
\definecolor{halfgray}{gray}{0.55}
\definecolor{ipython_frame}{RGB}{207, 207, 207}
\definecolor{ipython_bg}{RGB}{247, 247, 247}
\definecolor{ipython_red}{RGB}{186, 33, 33}
\definecolor{ipython_green}{RGB}{0, 128, 0}
\definecolor{ipython_cyan}{RGB}{64, 128, 128}
\definecolor{ipython_purple}{RGB}{170, 34, 255}
\definecolor{linkcolor}{HTML}{0000FF} % цвет гиперссылок
\definecolor{urlcolor}{HTML}{800080} % цвет ссылок
\definecolor{backcolour}{rgb}{0.95,0.95,0.92}

\usepackage{longtable}
\usepackage{wrapfig}
\usepackage{ragged2e}
\usepackage[nooneline]{caption}
\DeclareCaptionTextFormat{center}{\centering{#1}}
\captionsetup[table]{justification=raggedleft,
labelformat=empty,
labelsep=endash,
textformat=center,
position=top,
skip=5pt
}

\hyphenpenalty=10000% i don’t think hyphenation in presentations is a good idea, feel free to change however you like

\title{\LARGE{Системный анализ процессов химической технологии}}
\subtitle{\textcolor{tpugreen}{\textbf{Лабораторная работа №5}} \\ \textbf{Массивы NumPy}}
\author[]{Вячеслав Алексеевич Чузлов, \\
к.т.н., доцент ОХИ ИШПР}
\date{\today}

\begin{document}

\newcommand{\pythoninline}[1]{%
	\colorbox{white}{%
		\parbox[b][.6em]{\widthof{\mintinline[fontsize=\tiny]{ipython}{#1}}}{\mintinline[fontsize=\tiny]{ipython}{#1}}%
	}%
}

% notice usage of \titleframe and several other unconventional functions
% the reason being is custom backgrounds on these slides

\titleframe% title

%\tocframe{}% this custom frame accepts options for ToC



\section{Преобразование списков Python в массивы NumPy}
\begin{frame}[fragile]{Преобразование списков Python в массивы NumPy}
\scriptsize
\begin{itemize}
	\item Для того, чтобы создать объект массива NumPy из объекта списка Python, можно использовать функцию \texttt{np.array}:
\begin{minted}{pycon}
>>> import numpy as np
>>> np.array([1, 3, 5, 4, 2]) # Массив целочисленных значений
array([1, 3, 5, 4, 2])
\end{minted}

\item В отличие от стандартных списков Python, массивы NumPy могут содержать элементы только одного типа. Если типы элементов не совпадают, NumPy сделает попытку повышающего приведения типов:

\begin{minted}[firstnumber=last]{pycon}
>>> np.array([3.14, 4, 2, 3, 2.71])
array([3.14, 4.  , 2.  , 3.  , 2.71])
\end{minted}

\item В тех случаях, когда требуется явно задать тип результирующего массива, необходимо воспользоваться ключевым аргументом \texttt{dtype}:

\begin{minted}[firstnumber=last]{pycon}
>>> np.array([1, 3, 5, 4, 2], dtype='float32')
array([1., 3., 5., 4., 2.], dtype=float32)
\end{minted}
\end{itemize}
\vfill
\end{frame}

\begin{frame}[fragile]{Преобразование списков Python в массивы NumPy}
\scriptsize
\begin{itemize}
	\item В отличие от списков, массивы NumPy можно явным образом описать как многомерные:

\begin{minted}[firstnumber=last]{pycon}
>>> np.array([range(i, i + 3) for i in [2, 4, 6]])
array([[2, 3, 4],
       [4, 5, 6],
       [6, 7, 8]])
\end{minted}
\end{itemize}
\vfill
\end{frame}

\section{Создание массивов}
\begin{frame}[fragile]{Создание массивов}
\scriptsize
Массивы больших размеров эффективнее генерировать с помощью встроенных методов.
\begin{itemize}
\item Создаем массив целых чисел длины 10, заполненный нулями:
\begin{minted}{pycon}
>>> np.zeros(10, dtype=int)
array([0, 0, 0, 0, 0, 0, 0, 0, 0, 0])
\end{minted}
\item Создадим массив размером $3 \times 5$ значений с плавающей точкой, заполненный единицами:
\begin{minted}[firstnumber=last]{pycon}
>>> np.ones((3, 5), dtype=float)
array([[1., 1., 1., 1., 1.],
       [1., 1., 1., 1., 1.],
       [1., 1., 1., 1., 1.]])
\end{minted}
\item Создадим массив размером $3 \times 5$, заполненный значением $2.98$:
\begin{minted}[firstnumber=last]{pycon}
>>> np.full((3, 5), 2.98)
array([[2.98, 2.98, 2.98, 2.98, 2.98],
       [2.98, 2.98, 2.98, 2.98, 2.98],
       [2.98, 2.98, 2.98, 2.98, 2.98]])
\end{minted}
\end{itemize}
\vfill
\end{frame}

\begin{frame}[fragile]{Создание массивов}
\scriptsize
\begin{itemize}
\item Создадим массив, заполненный линейной последовательностью, начинающейся с $0$ и заканчивающейся $20$, с шагом $2$ (аналогично встроенной функции \mintinline{python}|range()|):
\begin{minted}[firstnumber=last]{pycon}
>>> np.arange(0, 20, 2)
array([ 0,  2,  4,  6,  8, 10, 12, 14, 16, 18])
\end{minted}
\item Создадим массив из пяти значений, равномерно располагающихся между $0$ и $1$:
\begin{minted}[firstnumber=last]{pycon}
>>> np.linspace(0, 1, 5)
array([0.  , 0.25, 0.5 , 0.75, 1.  ])
\end{minted}
\item Создадим массив размером $3 \times 3$ равномерно распределенных случайных значений от $0$ до $1$
\begin{minted}[firstnumber=last]{pycon}
>>> np.random.random((3, 3))
array([[0.28209864, 0.0014687 , 0.71104325],
       [0.58773349, 0.89099475, 0.87835838],
       [0.81093666, 0.70139097, 0.12816416]])
\end{minted}
\end{itemize}
\end{frame}

\begin{frame}[fragile]{Создание массивов}
\scriptsize
\begin{itemize}
	\item Создадим массив размером $3 \times 3$ случайных целых чисел в интервале $\left[0, 10\right]$
\begin{minted}[firstnumber=last]{pycon}
>>> np.random.randint(0, 10, (3, 3))
array([[0, 4, 0],
       [6, 7, 2],
       [0, 7, 3]])
\end{minted}
\end{itemize}
\vfill
\end{frame}

\section{Доступ к элементам массива}
\begin{frame}[fragile]{Доступ к элементам массива}
\scriptsize
Индексация массивов NumPy во многом похожа на индексацию списков в Python.
В одномерном массиве обратиться к $\mathrm{i}$-му (считая с $0$) значению можно по требуемому индексу в квадратных скобках, по аналогии со стандартными списками:

\begin{minted}{pycon}
>>> x1 = np.array([5, 0, 3, 3, 7, 9])
>>> x1[0]
5
>>> x1[4]
7
\end{minted}

Для индексации элементов с конца массива следует использовать отрицательные индексы:

\begin{minted}[firstnumber=last]{pycon}
>>> x1[-1]
9
>>> x1[-2]
7
\end{minted}
\end{frame}

\begin{frame}[fragile]{Доступ к элементам массива}
\scriptsize
Для обращения к элементам многомерного массива нужно указать кортеж индексов, разделенных запятыми:
\begin{minted}[firstnumber=last]{pycon}
>>> x2 = np.array([[3, 5, 2, 4], [7, 6, 8, 8], [1, 6, 7, 7]])
>>> x2
array([[3, 5, 2, 4],
       [7, 6, 8, 8],
       [1, 6, 7, 7]])
>>> x2[0, 0]
3
>>> x2[2, 0]
1
>>> x2[2, -1]
7
\end{minted}
\end{frame}

\begin{frame}[fragile]{Доступ к элементам массива}
\scriptsize
При помощи любой из указанных выше нотаций можно изменять значения элементов массива:

\begin{minted}[firstnumber=last]{pycon}
>>> x2[0, 0] = 24
>>> x2
array([[24,  5,  2,  4],
       [ 7,  6,  8,  8],
       [ 1,  6,  7,  7]])
\end{minted}

Следует помнить, что, в отличие от списков, массивы NumPy имеют фиксированный тип данных.
Если вставить в массив целых чисел значение с плавающей точкой, оно будет неявно усечено:

\begin{minted}[firstnumber=last]{pycon}
>>> x1[0] = 2.71828  # Это значение будет усечено!
>>> x1
array([2, 0, 3, 3, 7, 9])
\end{minted}
\vfill
\end{frame}

\subsection{Срезы массивов: доступ к подмассивам}
\begin{frame}[fragile]{Срезы массивов: доступ к подмассивам}
\scriptsize
\begin{itemize}
\item По аналогии с доступом к определенным элементам массива квадратные скобки используются для доступа к отдельным частям массива при помощи операции среза (\texttt{slicing}), обозначаемой знаком двоеточия (\texttt{:}).
\item Синтаксически срезы массивов NumPy соответствуют  срезам стандартных списков Python:
\\

\center{\texttt{x[начало:конец:шаг]}}
\\
\noindent при этом любое из значений можно не указывать, тогда по умолчанию будут приняты следующие значения: \texttt{начало = 0}, \texttt{конец = размер соответствующего измерения}, \texttt{шаг = 1}.
\end{itemize}
\subsubsection{Одномерные подмассивы}
\begin{alertblock}{\textbf{Одномерные подмассивы}}
\begin{minted}{pycon}
>>> x = np.arange(10)
>>> x
array([0, 1, 2, 3, 4, 5, 6, 7, 8, 9])
>>> x[:5]    # первые пять элементов
array([0, 1, 2, 3, 4])
>>> x[5:]    # элементы после индекса = 5
array([5, 6, 7, 8, 9])
>>> x[4:7]   # подмассив из середины
array([4, 5, 6])
\end{minted}
\end{alertblock}
\vfil
\end{frame}

\begin{frame}[fragile]{Срезы массивов: доступ к подмассивам}
\scriptsize
\begin{minted}[firstnumber=last]{pycon}
>>> x[::2]   # каждый второй элемент
array([0, 2, 4, 6, 8])
>>> x[1::2]  # каждый второй элемент, начиная с индекса 1
array([1, 3, 5, 7, 9])
\end{minted}
\begin{itemize}
\item Также сохранена возможность использования отрицательного значения параметра \texttt{\textbf{шаг}}. Тогда значения по умолчанию для параметров \texttt{\textbf{начало}} и \texttt{\textbf{конец}} будут поменяны местами. Это быстрый способ перевернуть массив:
\end{itemize}

\begin{minted}[firstnumber=last]{pycon}
>>> x[::-1]    # все элементы в обратном порядке
array([9, 8, 7, 6, 5, 4, 3, 2, 1, 0])
>>> x[5::-1]   # все элементы в обратном порядке, начиная с индекса 5
array([5, 4, 3, 2, 1, 0])
>>> x[5::-2]  # каждый второй элемент в обратном порядке, начиная с индекса 5
array([5, 3, 1])
\end{minted}
\vfill
\end{frame}

\begin{frame}[fragile]{Срезы массивов: доступ к подмассивам}
\scriptsize
\subsubsection{Многомерные подмассивы}
\begin{alertblock}{\textbf{Многомерные подмассивы}}
\begin{minted}{pycon}
>>> x2
array([[24,  5,  2,  4],
       [ 7,  6,  8,  8],
       [ 1,  6,  7,  7]])
>>> x2[:2, :3]             # две строки, три столбца
array([[24,  5,  2],
       [ 7,  6,  8]])
>>> x2[:3, ::2]            # все строки, каждый второй столбец
array([[24,  2],
       [ 7,  8],
       [ 1,  7]])
>>> x2[::-1, ::-1]         # обратный порядок строк и столбцов
array([[ 7,  7,  6,  1],
       [ 8,  8,  6,  7],
       [ 4,  2,  5, 24]])
\end{minted}
\end{alertblock}
\vfil
\end{frame}

\begin{frame}[fragile]{Срезы массивов: доступ к подмассивам}
\scriptsize
\subsubsection{Доступ к строкам и столбцам массива}
\begin{alertblock}{\textbf{Доступ к строкам и столбцам массива}}
Распространенной задачей является доступ к отдельным строкам или столбцам массива.
Получить такой доступ можно при помощи комбинации операций индексации и среза:
\begin{minted}[firstnumber=last]{pycon}
>>> print(x2[:, 0])   # первый столбец массива x2
[24  7  1]
>>> print(x2[0, :])   # первая строка массива x2
[24  5  2  4]
\end{minted}
При необходимости получения доступа к строке, операция взятия среза может быть опущена для более лаконичной записи:
\begin{minted}{pycon}
>>> print(x2[0])  # эквивалентно x2[0, :]
[24  5  2  4]
\end{minted}
\end{alertblock}
\vfill
\end{frame}

\begin{frame}[fragile]{Срезы массивов создают разделяемые ссылки}
\scriptsize
\begin{itemize}
\item Срезы массивов возвращают \textcolor{extraorange}{\textbf{разделяемые ссылки}} (синонимы), а не \textcolor{extraorange}{\textbf{копии}} данных массива.
\item Этим срезы массивов библиотеки NumPy отличаются от срезов списков языка Python (срезы списков создают новые объекты):
\begin{minted}{pycon}
>>> print(x2)
[[24  5  2  4]
 [ 7  6  8  8]
 [ 1  6  7  7]]
\end{minted}

Получим из него матрицу $2 \times 2$:

\begin{minted}[firstnumber=last]{pycon}
>>> x2_sub = x2[:2, :2]
>>> print(x2_sub)
[[24  5]
 [ 7  6]]
\end{minted}
\end{itemize}
\vfill
\end{frame}

\begin{frame}[fragile]{Срезы массивов создают разделяемые ссылки}
\scriptsize
Теперь, если изменить значения этой матрицы, исходный массив также поменялся:

\begin{minted}[firstnumber=last]{pycon}
>>> x2_sub[0, 0] = 100
>>> print(x2_sub)
[[100   5]
 [  7   6]]
>>> print(x2)
[[100   5   2   4]
 [  7   6   8   8]
 [  1   6   7   7]]
\end{minted}
\vfill
\end{frame}

\subsection{Создание копий массивов}
\begin{frame}[fragile]{Создание копий массивов}
\scriptsize
В ряде случаев требуется явно скопировать содержимое массива или его части. Для решения данной задачи существует метод \mintinline{python}|copy()|:

\begin{minted}[firstnumber=last]{pycon}
>>> x2_sub_copy = x2[:2, :2].copy()
>>> print(x2_sub_copy)
[[100   5]
 [  7   6]]
\end{minted}

Теперь, если изменить значения этого подмассива, то исходный массив не изменится:

\begin{minted}[firstnumber=last]{pycon}
>>> x2_sub_copy[0, 0] = 24
>>> print(x2_sub_copy)
[[24  5]
 [ 7  6]]
>>> print(x2)
[[100   5   2   4]
 [  7   6   8   8]
 [  1   6   7   7]]
\end{minted}
\vfill
\end{frame}

\section{Вычисления с массивами NumPy}
\sectionframe

\subsection{Арифметические функции с массивами}
\begin{frame}[fragile]{Арифметические функции с массивами}
\scriptsize
Универсальные функции NumPy могут быть легко использованы, т.к. основаны на нативных арифметических операторах Python. Доступны обычные операторы сложения, вычитания, умножения и деления:

\begin{minted}{pycon}
>>> x = np.arange(4)
>>> print('x      =', x)
x      = [0 1 2 3]
>>> print('x + 5  =', x + 5)
x + 5  = [5 6 7 8]
>>> print('x - 5  =', x - 5)
x - 5  = [-5 -4 -3 -2]
>>> print('x * 2  =', x * 2)
x * 2  = [0 2 4 6]
>>> print('x / 2  =', x / 2)
x / 2  = [0.  0.5 1.  1.5]
>>> print('x // 2 =', x // 2)
x // 2 = [0 0 1 1]
\end{minted}
\vfill
\end{frame}

\begin{frame}[fragile]{Арифметические функции с массивами}
\scriptsize
Определены также унарная универсальная функция изменения знака, оператор \mintinline{python}|**| для возведения в степень и оператор \mintinline{python}|%| для деления по модулю:
\begin{minted}[firstnumber=last]{pycon}
>>> print('-x     =', -x)
-x     = [ 0 -1 -2 -3]
>>> print('x ** 2 =', x ** 2)
x ** 2 = [0 1 4 9]
>>> print('x % 2  =', x % 2)
x % 2  = [0 1 0 1]
\end{minted}
\vfill
\end{frame}

\begin{frame}[fragile]{Арифметические функции с массивами}
\scriptsize
Данные операции могут быть использованы в выражениях любыми способами с соблюдением стандартных приоритетов:

\begin{minted}[firstnumber=last]{pycon}
>>> -(0.5 * x + 1) ** 2
array([-1.  , -2.25, -4.  , -6.25])
\end{minted}

Все арифметические операторы~-- удобные аналоги для встроенных функций библиотеки NumPy. Например, оператор \mintinline{python}|+| является аналогом функции \mintinline{python}|add()|:

%\begin{minted}[firstnumber=last]{pycon}
%>>> np.add(x, 5)
%array([5, 6, 7, 8])
%\end{minted}
\begin{table}[h!]
%\caption{Реализованные в бибилотеке NumPy арифметические операторы}
\label{tab:арифм.опер.np}
\begin{tabular*}{\linewidth}{p{0.11\linewidth}p{0.28\linewidth}p{0.54\linewidth}}
\hline
{Оператор} & {Эквивалентная функция} & {Описание} \\
\hline

\mintinline{ipython}|+| & \mintinline{ipython}|np.add| & Сложение: \mintinline{ipython}|1 + 1 = 2| \\

\mintinline{ipython}|-| & \mintinline{ipython}|np.subtract| & Вычитание: \mintinline{ipython}|3 - 2 = 1| \\

\mintinline{ipython}|-| & \mintinline{ipython}|np.negative| & Унарная операция изменения знака: \mintinline{ipython}|-2| \\

\mintinline{ipython}|*| & \mintinline{ipython}|np.multiply| & Умножение: \mintinline{ipython}|2 * 3 = 6| \\

\mintinline{ipython}|/| & \mintinline{ipython}|np.divide| & Деление: \mintinline{ipython}|3 / 2 = 1.5| \\

\mintinline{ipython}|//| & \mintinline{ipython}|np.floor_divide| & Деление с округлением в меньшую сторону: \mintinline{ipython}|3 // 2 = 1| \\

\mintinline{ipython}|**| & \mintinline{ipython}|np.power| & Возведение в степень: \mintinline{ipython}|3 ** 2 = 9| \\

\mintinline{ipython}|%| & \mintinline{ipython}|np.mod| & Модуль/остаток: \mintinline{ipython}|5 % 2 = 1| \\
\hline
\end{tabular*}
\end{table}
\vfil
\end{frame}

\subsection{Получение абсолютного значения}
\begin{frame}[fragile]{Получение абсолютного значения}
\scriptsize
Наряду со встроенными арифметическими операторами, с массивами NumPy можно использовать стандартную функцию \mintinline{python}|abs()| языка Python для получения абсолютного значения:

\begin{minted}{pycon}
>>> x = np.array([-2, -1, 0, 1, 2])
>>> abs(x)
array([2, 1, 0, 1, 2])
\end{minted}

Аналогичная универсальная функция NumPy~-- \mintinline{python}|np.absolute()|, доступна также под псевдонимом \mintinline{python}|np.abs()|:

\begin{minted}[firstnumber=last]{pycon}
>>> np.absolute(x)
array([2, 1, 0, 1, 2])
>>> np.abs(x)
array([2, 1, 0, 1, 2])
\end{minted}
\vfil
\end{frame}

\subsection{Тригонометрические функции}
\begin{frame}[fragile]{Тригонометрические функции}
\scriptsize
Библиотека NumPy предоставляет набор тригонометрических функций:

\begin{minted}{pycon}
>>> alpha = np.linspace(0, np.pi, 3)
>>> print('alpha      = ', alpha)
alpha      =  [0.         1.57079633 3.14159265]
>>> print('sin(alpha) = ', np.sin(alpha))
sin(alpha) =  [0.0000000e+00 1.0000000e+00 1.2246468e-16]
>>> print('cos(alpha) = ', np.cos(alpha))
cos(alpha) =  [ 1.000000e+00  6.123234e-17 -1.000000e+00]
>>> print('tan(alpha) = ', np.tan(alpha))
tan(alpha) =  [ 0.00000000e+00  1.63312394e+16 -1.22464680e-16]
\end{minted}

Значения вычисляются в пределах точности конкретной вычислительной машины, вследствие чего некоторые из них не всегда точно равны нулю, хотя должны.
\vfill
\end{frame}

\begin{frame}[fragile]{Тригонометрические функции}
\scriptsize
Определены также  и обратные тригонометрические функции:

\begin{minted}[firstnumber=last]{pycon}
>>> x = [-1, 0, 1]
>>> print('x = ', x)
x =  [-1, 0, 1]
>>> x = [-1, 0, 1]
>>> print('x = ', x)
x =  [-1, 0, 1]
>>> print('arcsin(x) = ', np.arcsin(x))
arcsin(x) =  [-1.57079633  0.          1.57079633]
>>> print('arccos(x) = ', np.arccos(x))
arccos(x) =  [3.14159265 1.57079633 0.        ]
>>> print('arctan(x) = ', np.arctan(x))
arctan(x) =  [-0.78539816  0.          0.78539816]
\end{minted}
\vfill
\end{frame}

\subsection{Показательные функции и логарифмы}
\begin{frame}[fragile]{Показательные функции и логарифмы}
\scriptsize
Показательные функции~-- один из  распространенных типов операций, доступных в NumPy:

\begin{minted}{pycon}
>>> x = [1, 2, 3]
>>> print('x   = ', x)
x   =  [1, 2, 3]
>>> print('e^x = ', np.exp(x))
e^x =  [ 2.71828183  7.3890561  20.08553692]
>>> print('2^x = ', np.exp2(x))
2^x =  [2. 4. 8.]
>>> print('3^x = ', np.power(3, x))
3^x =  [ 3  9 27]
\end{minted}
\vfill
\end{frame}

\begin{frame}[fragile]{Показательные функции и логарифмы}
\scriptsize
Определены также и логарифмы. Простейшая функция \mintinline{python}|np.log()| возвращает натуральный логарифм числа. Если Вам требуется логарифм по основанию $2$ или $10$, они также доступны:

\begin{minted}[firstnumber=last]{pycon}
>>> x = [1, 2, 4, 10]
>>> print('x        =', x)
x        = [1, 2, 4, 10]
>>> print('ln(x)    =', np.log(x))
ln(x)    = [0.         0.69314718 1.38629436 2.30258509]
>>> print('log2(x)  =', np.log2(x))
log2(x)  = [0.         1.         2.         3.32192809]
>>> print('log10(x) =', np.log10(x))
log10(x) = [0.         0.30103    0.60205999 1.        ]
\end{minted}
\vfill
\end{frame}


\section{Функции агрегирования}
\sectionframe

\subsection{Суммирование значений массива}
\begin{frame}[fragile]{Суммирование значений массива}
\scriptsize
В стандартном Python данная задача решается при помощи встроенной функции \mintinline{ipython}|sum()|. Синтаксис этой функции крайне похож на функцию \mintinline{ipython}|sum()| библиотеки NumPy:

\begin{minted}{pycon}
>>> import numpy as np
>>> arr = np.random.random(100)
>>> sum(arr)
49.496408327779065
>>> np.sum(arr)
49.49640832777906
\end{minted}

NumPy версия функции \mintinline{ipython}|sum()| работает намного быстрее:

\begin{minted}[firstnumber=last]{pycon}
>>> big_array = np.random.rand(1000000)
>>> %timeit sum(big_array)  # Магическая команда в IPython
62.8 ms |$\pm$| 598 |$\mu$|s per loop (mean |$\pm$| std. dev. of 7 runs, 10 loops each
>>> %timeit np.sum(big_array)
522 |$\mu$|s |$\pm$| 5.53 |$\mu$|s per loop (mean |$\pm$| std. dev. of 7 runs, 1000 loops each)
\end{minted}

Важно заметить, что функции \mintinline{ipython}|sum()| и \mintinline{ipython}|np.sum()| не идентичны. Так, смысл их опциональных аргументов отличается и функция \mintinline{ipython}|np.sum()| может работать с многомерными массивами.
\vfill
\end{frame}

\subsection{Минимум и максимум}
\begin{frame}[fragile]{Минимум и максимум}
\scriptsize
В стандартном Python определены встроенные функции \mintinline{ipython}|min()| и \mintinline{ipython}|max()|, служащие для вычисления минимального и максимального значений любой коллекции:

\begin{minted}{pycon}
>>> min(big_array), max(big_array)
(1.4200179176970806e-07, 0.9999998044567884)
\end{minted}

Соответствующие функций NumPy имеют аналогичный синтаксис и работают быстрее:

\begin{minted}[firstnumber=last]{pycon}
>>> np.min(big_array), np.max(big_array)
(1.4200179176970806e-07, 0.9999998044567884)
>>> %timeit min(big_array)
43.4 ms |$\pm$| 719 |$\mu$|s per loop (mean |$\pm$| std. dev. of 7 runs, 10 loops each)
>>> %timeit np.min(big_array)
304 |$\mu$|s |$\pm$| 2.03 |$\mu$|s per loop (mean |$\pm$| std. dev. of 7 runs, 1000 loops each)
\end{minted}

Существует возможность вычисления некоторых сводных показателей путем вызова соответствующих методов объекта массива NumPy для более лаконичной записи:

\begin{minted}[firstnumber=last]{pycon}
>>> print(big_array.min(), big_array.max(), big_array.sum())
1.4200179176970806e-07 0.9999998044567884 499940.4506893933
\end{minted}
\vfill
\end{frame}

\section{Примеры}
\sectionframe


\begin{frame}[fragile]{\textcolor{tpugreen}{\textbf{Пример~1}}}
\scriptsize
Заполнить матрицу $\mathrm{a}(3, 3)$ случайными целыми числами от $1$ до $10$. Вывести значения ее элементов. Вычислить сумму элементов первой и последней строк данной матрицы и вывести результат вычисления.
\vfill
\begin{minted}{pycon}
>>> import numpy as np
>>> a = np.random.randint(1, 10, (3, 3))
>>> a
array([[5, 2, 7],
       [7, 2, 1],
       [8, 2, 3]])
>>> s = a[0].sum() + a[-1].sum()
>>> s
27
\end{minted}
\vfill
\end{frame}


\begin{frame}[fragile]{\textcolor{tpugreen}{\textbf{Пример~2}}}
\scriptsize
Матрица $p$ размером $2 \times 3$ заполнена случайными целыми числами в диапазоне $\left[-3; 3\right]$. Массив $x$ содержит $10$ случайных чисел в диапазоне $\left[-10; 10\right]$. Необходимо найти сумму минимального элемента первой строки матрицы $p$ и максимального элемента массива $x$ и поделить ее на количество элементов массива $x$, больших нуля.
\vfill
\begin{minted}[fontsize=\fontsize{8pt}{8pt}]{pycon}
>>> import numpy as np
>>> p = np.random.randint(-3, 3, (2, 3))
>>> p
array([[-2,  2,  2],
       [ 0, -2,  2]])
>>> x = np.random.randint(-10, 10, 10)
>>> x
array([ 7, -7, -8, -7, -8, -3, -8,  6, -5,  6])
>>> min_p = p[0].min()
>>> max_x = x.max()
>>> min_p, max_x
(-2, 7)
>>> xshape = x[x > 0].shape
>>> xshape
(3,)
>>> (max_x + min_p) / xshape[0]
1.6666666666666667
\end{minted}
\vfill
\end{frame}

\begin{frame}[fragile]{\textcolor{tpugreen}{\textbf{Пример~3}}}
\scriptsize
Заполнить матрицу $\mathrm{a}(4, 4)$ случайными целыми числами от $1$ до $100$. Вывести значения ее элементов. Найти максимальный элемент в каждой строке. Среди максимальных элементов каждой строки найти минимальный.
\vfill
\begin{minted}{pycon}
>>> import numpy as np
>>> a = np.random.randint(1, 100, (4, 4))
>>> a
array([[62, 97, 91, 19],
       [34, 34, 92, 70],
       [96, 92, 46, 80],
       [97,  2, 75, 52]])
>>> max_of_rows = a.max(axis=1)  # axis=0 - поиск максимума по столбцам
>>> max_of_rows
array([97, 92, 96, 97])
>>> max_of_rows.min()
92
\end{minted}
\vfill
\end{frame}

\begin{frame}[fragile]{\textcolor{tpugreen}{\textbf{Пример~4}}}
\scriptsize
Заполнить матрицу $\mathrm{a}(4, 4)$ случайными числами от $-3$ до $6$. Вывести значения ее элементов. Вычислить среднее арифметическое значений неотрицательных элементов каждого столбца данной матрицы.
\vfill
\begin{minted}{pycon}
>>> import numpy as np
>>> a = np.random.randint(-3, 6, (4, 4))
>>> a
array([[ 3, -3,  0, -1],
       [-2,  1, -1,  4],
       [ 3, -2,  0,  5],
       [-1, -2,  0,  0]])
>>> mean_of_cols = a.mean(axis=0, where=a>=0)
>>> mean_of_cols
array([3., 1., 0., 3.])
\end{minted}
\vfill
\end{frame}


\section{Задания}
\sectionframe

\begin{frame}[fragile]{\textcolor{tpugreen}{\textbf{Задание~1}}}
\scriptsize
Заполнить матрицу $\mathrm{a}\left(3, 3\right)$ случайными целыми числами от $-5$ до $5$. Вывести значения ее элементов на экран. Найти произведение минимального элемента матрицы на сумму ее положительных элементов. Вывести результат.
\vfill
\end{frame}

\begin{frame}[fragile]{\textcolor{tpugreen}{\textbf{Задание~2}}}
\scriptsize
Заполнить матрицу $\mathrm{a}\left(5, 10\right)$ случайными целыми числами от $0$ до $9$. Вывести значения элементов матрицы. Найти столбец матрицы с максимальной суммой элементов и вывести эту сумму на экран.
\vfill
\end{frame}

\begin{frame}[fragile]{\textcolor{tpugreen}{\textbf{Задание~3}}}
\scriptsize
Найти сумму положительных ($p$) и количество отрицательных ($o$) элементов массива $\mathrm{z}\left(10\right)$, заполненного случайными целыми числами в диапазоне $\left[-7; 5\right]$ и минимальный элемент $\left(b_{min}\right)$ второго столбца матрицы $\mathrm{b\left(5, 5\right)}$, заполненной случайными числами в диапазоне $\left[-10; 20\right]$, и вычислить значения элементов массива $x_i$:
\vfill
\begin{normalsize}
\begin{equation*}
	x_i = \dfrac{\sqrt{p}}{o + a \cdot c} + b_{min} \cdot k^2 + z_i
\end{equation*}
\end{normalsize}
\vfill
где $a = 2.5\mathrm{e-}3; \quad c = 175; \quad k = 8$.
\vfill
\end{frame}

\contactsframe[\Large \textbf{Благодарю за внимание!}]{

\includegraphics[width=.05\textwidth]{pics/home} \quad Учебный корпус №2, ауд. 136 \\
\includegraphics[width=.05\textwidth]{pics/mail} \quad chuva@tpu.ru \\
\includegraphics[width=.03\textwidth]{pics/tel} \quad +7-962-782-66-15
}

\end{document}

